\documentclass[12pt]{article}
\usepackage[utf8]{inputenc}
\usepackage[T1]{fontenc}
\usepackage[margin=1in]{geometry}
\usepackage{setspace}
\usepackage{xcolor}
\usepackage[colorlinks=true, urlcolor=blue!60!black, linkcolor=blue!60!black]{hyperref}
\setlength{\parindent}{2em}
\setlength{\parskip}{1.25em}
\renewcommand{\baselinestretch}{1.0}
\begin{document}
\begin{center}
{\LARGE \textit{Buck v. Bell}, 274 U.S. 200 (1927)\par}
\vspace{0.8em}
{\large Case Brief --- Four-Page Edition\par}
\end{center}
\vspace{1em}

\section{Memory Jogger}

The Court upheld Virginia's compulsory eugenic sterilization statute against due process and equal protection challenges, analogizing forced salpingectomy to compulsory vaccination and declaring that ``[t]hree generations of imbeciles are enough'' (\href{https://supreme.justia.com/cases/federal/us/274/200/}{274 U.S. 200}).

\section{Detailed Case Facts}

Virginia's 1924 statute authorized sterilizing ``mental defectives'' in state institutions, reciting that heredity played an important role in transmitting insanity and imbecility and that sterilized inmates could be safely discharged. It required a petition by the superintendent, notice to the inmate and a court-appointed guardian, a hearing, and appeals to the county circuit court and the Virginia Supreme Court of Appeals.

Carrie Buck, eighteen, had been committed to the Virginia State Colony for Epileptics and Feeble-Minded. The record described her as a ``feeble minded white woman,'' the daughter of a ``feeble minded mother,'' and the mother of an ``illegitimate feeble minded child.'' Later research by Paul Lombardo showed that Carrie had become pregnant after being raped by her foster parents' nephew and was committed to conceal the scandal; neither she nor her daughter Vivian was intellectually disabled. Superintendent Priddy selected Buck as a test case, and her counsel, a former colony board member, called no witnesses and conceded her ``feeblemindedness.''

\section{Procedural History}

The colony's board approved sterilization in September 1924. Buck's guardian appealed to the Amherst County Circuit Court, which upheld the order, finding Buck ``the probable potential parent of socially inadequate offspring, likewise afflicted.'' After Priddy died, Dr.\ John Bell was substituted. The Virginia Supreme Court of Appeals affirmed, and the U.S. Supreme Court affirmed on May 2, 1927. Bell sterilized Buck that October.

\section{Judicial Votes}

Affirmed, \textbf{8--1}.
\begin{itemize}
\item \textbf{Majority:} Justice Holmes, joined by Chief Justice Taft and Justices Van Devanter, McReynolds, Brandeis, Sutherland, Sanford, and Stone.
\item \textbf{Dissent:} Justice Butler, without opinion.
\end{itemize}

\section{Holding}

A state statute authorizing compulsory sterilization of inmates of state institutions adjudicated, after thorough procedural safeguards, to carry hereditary insanity or imbecility does not violate due process; it falls within the police power. Limiting the law to institutional inmates does not violate equal protection, because a state may address a problem incrementally as its means allow. Holmes's rhetoric about ``three generations'' did not set a doctrinal threshold, and the Court did not address sterilizing non-institutionalized persons or sterilization as punishment.

\section{Analysis of Opinions}

\textbf{Majority (Holmes).} In barely three pages, Holmes first found procedural due process satisfied: after the statute's notice, guardian, hearing, and appeal provisions, ``the rights of the patient are most carefully considered.'' On substance, he deferred entirely to the legislature's eugenic premises and the finding that Buck would bear ``socially inadequate offspring.'' Analogizing to wartime conscription, he reasoned that if the state can demand the lives of its ``best citizens,'' it may demand ``these lesser sacrifices'' from those ``who already sap the strength of the State,'' and he invoked \textit{Jacobson v. Massachusetts}. He dismissed equal protection as ``the usual last resort of constitutional arguments,'' adding that sterilization freed asylum beds and so brought more people within the law's reach.

\textbf{Dissent (Butler).} Butler dissented without opinion. Historians often attribute his vote to his Catholicism and commitment to individual liberty, but he left no doctrinal counter-framework.

\section{Examples: Future Applications}

\begin{enumerate}
\item \textbf{Same-Side: Institutional Vasectomy Statute.} A state copies Virginia's law for male psychiatric patients found, after a hearing with counsel and appeal, to carry hereditary severe impairment. Robust procedure and uniform institutional application fall within \textit{Buck}.
\item \textbf{Same-Side: Incremental Vaccination Mandate.} A state mandates vaccination only in state facilities and public dormitories. Under \textit{Buck}, a state may act ``within the lines'' as its means allow.
\item \textbf{Opposite-Side: Summary Sterilization.} A superintendent orders sterilization based on a nurse's notes, with no notice, hearing, or review. Holmes relied on Virginia's careful procedures, so this fails \textit{Buck}'s procedural floor.
\item \textbf{Opposite-Side: Punitive Sterilization of Thieves.} A statute sterilizes three-time larcenists but exempts embezzlers. As \textit{Skinner v. Oklahoma} (1942) held, arbitrary punitive sterilization violates equal protection.
\item \textbf{Fence-Sitter: Reversible Contraception for Welfare Recipients.} After full hearings, a state requires reversible implants for community-dwelling parents with hereditary intellectual disabilities who receive assistance. The measure is less permanent and procedurally careful, but it reaches free citizens outside state custody.
\end{enumerate}

\section{Critique}

\textbf{Anti-canon.} \textit{Buck} sits in the anti-canon with \textit{Dred Scott}, \textit{Plessy}, and \textit{Korematsu}. The \textit{Jacobson} analogy is a false equivalence between a temporary vaccination against epidemic disease and permanent destruction of reproductive capacity; Holmes accepted a sham record from collusive litigation; and his equal protection reasoning reduced the clause to a rubber stamp. The decision catalyzed over 60,000 sterilizations, was cited by Nazi defendants at Nuremberg, and has never been formally overruled.

\textbf{Progressive and originalist views.} Feminist and disability-rights scholars see class bias, gender policing, and ableism masquerading as science, and note that Holmes and Brandeis, heroes of judicial restraint, showed how uncritical deference invites abuse. Originalists treat bodily integrity as a deeply rooted common-law right, and Justice Thomas in \textit{Box v. Planned Parenthood} (2019) cited \textit{Buck} as a cautionary monument to Progressive social engineering.

\section{Key Quotations}

\begin{quote}
``It is better for all the world if, instead of waiting to execute degenerate offspring for crime or to let them starve for their imbecility, society can prevent those who are manifestly unfit from continuing their kind. The principle that sustains compulsory vaccination is broad enough to cover cutting the Fallopian tubes. Jacobson v. Massachusetts, 197 U. S. 11. Three generations of imbeciles are enough.''
\end{quote}
\noindent --- Justice Holmes, Opinion of the Court, 274 U.S. at 207.

\begin{quote}
``We have seen more than once that the public welfare may call upon the best citizens for their lives. It would be strange if it could not call upon those who already sap the strength of the State for these lesser sacrifices, often not felt to be such by those concerned, in order to prevent our being swamped with incompetence.''
\end{quote}
\noindent --- Justice Holmes, Opinion of the Court, 274 U.S. at 207.

\begin{quote}
``It is the usual last resort of constitutional arguments to point out shortcomings of this sort. But the answer is that the law does all that is needed when it does all that it can, indicates a policy, applies it to all within the lines, and seeks to bring within the lines all similarly situated so far and so fast as its means allow.''
\end{quote}
\noindent --- Justice Holmes, Opinion of the Court, 274 U.S. at 208.

\end{document}
